\section{Analysis and discussion}
\label{sec:analysis}

\subsection{Why the inversion is below random rather than at chance}
\label{sec:analysis:inversion}

The below-random orientation of test $\auroc$ under frozen natural-image
backbones is not noise but a structured failure: the linear predictor learns
a direction $w \in \mathbb{R}^{d}$ that encodes ``stable under
$\{\strongfam, \strongholdone\}$'' as small $\langle w, z_{t} - z_{c} \rangle$,
and this direction extrapolates with the \emph{opposite} sign when the unseen
family \strongholdtwo\ rotates the latent geometry the other way. Because the
third family's nuisance is operation-disjoint from the training families, the
rotation does not partially align --- it anti-aligns. Figure~\ref{fig:pair-score-histograms}
shows this structure directly: under frozen OpenAI CLIP, the test-split
stable-pair mean cosine distance ($0.281$) sits \emph{above} the
changing-pair mean ($0.233$), a $-0.049$ inversion gap. Under
BiomedCLIP, the gap narrows to $-0.034$. Under DermLIP, the ordering
flips to the correct direction with a $+0.111$ gap (stable
$0.236$, changing $0.346$). The natural-image inversion is large enough
relative to the within-class spread that $\auroc$ sits in the
$0.25$--$0.29$ band rather than at chance.

\subsection{Why DermLIP transfers under this scaffold and BiomedCLIP doesn't}
\label{sec:analysis:dermlip-vs-biomedclip}

The $+0.66$ $\auroc$ swing from EXP-006b (OpenAI CLIP) to EXP-007 (DermLIP)
cannot be explained by architecture, scaffold, or optimiser: all three are
held byte-identical. The plausible mechanistic accounts are
\begin{enumerate}[label=(\roman*)]
\item DermLIP's CLIP loss over $\sim$1\,M dermatology image-text pairs has
      organised the embedding space such that a single linear direction
      captures stable-pair similarity across all three nuisance families
      uniformly, including \strongholdtwo; or
\item DermLIP saw the HAM10000 raw images during pretraining and its
      embedding space is well-fit to HAM10000 lesion identity in a way that
      survives the synthetic post-hoc nuisance.
\end{enumerate}
EXP-008 rules out a third candidate --- ``any medical-image pretraining'' ---
by showing that BiomedCLIP, trained on $\sim$15\,M PMC-15M figure-caption
pairs without HAM10000, lifts test $\auroc$ by only $+0.04$ over web CLIP.
EXP-008 cannot, however, choose between (i) and (ii). EXP-009 (a
non-HAM10000 dermoscopy SSL pretrain, \S\ref{sec:followup}) is the
experiment that would.

\subsection{Loss versus AUROC}
\label{sec:analysis:loss-vs-auroc}

MSE loss delta does not predict train $\auroc$ across optimisers. From
the per-run training reports
(\texttt{<run\_id>/artifacts/reports/jepa\_training\_report.json} on the
run archive), EXP-006a (Adam, MLP) reduces train MSE by $\sim$$3.3\%$
and reaches train $\auroc = 0.893$, while EXP-004 (SGD, linear) reduces
train MSE by $\sim$$13\%$ and reaches train $\auroc = 0.900$. The
optimisers move the predictor's function very differently per unit of
L2-loss reduction in $768$-d normalised space. Reporting both metrics is
essential, and headline-$\auroc$ comparisons should not be replaced by
loss-only comparisons.

\subsection{What is not yet claimed}
\label{sec:analysis:non-claim}

We do not claim that frozen DermLIP ``solves'' the proxy in a transfer
sense. We claim that frozen DermLIP plus a linear scaffold reaches
$\auroc = 0.944 \pm 0.003$ on \strongholdtwo\ under the specific evaluation
protocol of \S\ref{sec:method}, with HAM10000 contamination unpartitioned.
The result is consistent with two hypotheses --- dermoscopy-domain transfer
and HAM10000 image-level overlap during DermLIP pretraining --- and EXP-009
is required to choose between them. The BiomedCLIP partition rules out only
the broader ``any medical-image pretraining is sufficient'' alternative.
